\subsection{Inner regular set functions}

DEFINITION 1. --- Let T be a topological space, and let \(\mathfrak{B}(T)\) be the Borel tribe of T; let I be a mapping of \(\mathfrak{B}(T)\) into \(\overline{\mathbf{R}}_{+}\) .

\begin{enumerate}
\def\labelenumi{\alph{enumi})}
\tightlist
\item
  I is said to be countably additive if, for every sequence \((\mathbf{A}_{n})\) of pairwise disjoint elements of \(\mathfrak{B}(\mathbf{T})\) ,
\end{enumerate}

\[
\mathrm{I} \left(\bigcup_ {n} \mathrm{A} _ {n}\right) = \sum_ {n} \mathrm{I} (\mathrm{A} _ {n}).\tag{4}
\]

\begin{enumerate}
\def\labelenumi{\alph{enumi})}
\setcounter{enumi}{1}
\tightlist
\item
  I is said to be inner regular if, for every set \(\mathbf{A} \in \mathfrak{B}(\mathbf{T})\) ,
\end{enumerate}

\[
\mathrm{I} (\mathrm{A}) = \sup _ {\mathrm{K}} \mathrm{I} (\mathrm{K}),\tag{5}
\]

where K runs over the set of compact subsets of A.

\begin{enumerate}
\def\labelenumi{\alph{enumi})}
\setcounter{enumi}{2}
\tightlist
\item
  I is said to be bounded (resp. locally bounded) if \(\mathrm{I(T)} < +\infty\) (resp. if every point \(x \in \mathbf{T}\) admits an open neighborhood V such that \(\mathrm{I(V)} < +\infty\) ).
\end{enumerate}

Remarks. --- 1) The condition a) clearly implies that I is an increasing mapping of \(\mathfrak{B}(\mathbf{T})\) (ordered by inclusion) into \(\overline{\mathbf{R}}_{+}\) .

\begin{enumerate}
\def\labelenumi{\arabic{enumi})}
\setcounter{enumi}{1}
\item
  Suppose that I is countably additive; let \((\mathbf{A}_n)_{n\in \mathbf{N}}\) be an increasing sequence of Borel sets, and let \(\mathbf{A} = \bigcup_{n\in \mathbf{N}}\mathbf{A}_n\) . The sets \(\mathrm{D}_0 = \mathrm{A}_0\) , \(\mathrm{D}_n = \mathrm{A}_n - \mathrm{A}_{n - 1}\) being pairwise disjoint, and their union being A, we have \(\mathrm{I}(\mathrm{A}) = \sum_{n} \mathrm{I}(\mathrm{D}_{n}) = \lim_{n \to \infty} \mathrm{I}(\mathrm{A}_{n})\) . Similarly, if \((\mathrm{B}_{n})\) is a decreasing sequence of Borel sets, and if \(\mathrm{I}(\mathrm{B}_{0}) < +\infty\) , then \(\mathrm{I}\left(\bigcap_{n} \mathrm{B}_{n}\right) = \lim_{n \to \infty} \mathrm{I}(\mathrm{B}_{n})\) : it suffices to apply the preceding to the sets \(A_{n} = B_{0} - B_{n}\) .
\item
  Let \((\mathbf{A}_n)\) be any sequence of Borel sets of T. If I is countably additive, then \(\mathrm{I}\left(\bigcup_{n} \mathrm{A}_n\right) \leqslant \sum_{n} \mathrm{I}(\mathrm{A}_n)\) . By the preceding remark, it suffices to establish this inequality for a finite sequence. One immediately reduces to the case of two sets \(A_{1}\) and \(A_{2}\) ; but the relation (4) implies that
\end{enumerate}

\[
\mathrm{I} \left(\mathrm{A} _ {1} \cup \mathrm{A} _ {2}\right) = \mathrm{I} \left(\mathrm{A} _ {1}\right) + \mathrm{I} \left(\mathrm{A} _ {2} - \left(\mathrm{A} _ {1} \cap \mathrm{A} _ {2}\right)\right) \leqslant \mathrm{I} \left(\mathrm{A} _ {1}\right) + \mathrm{I} \left(\mathrm{A} _ {2}\right).
\]

The desired inequality then follows immediately.

\begin{enumerate}
\def\labelenumi{\arabic{enumi})}
\setcounter{enumi}{3}
\item
  If I is a countably additive and locally bounded function, the preceding remark implies at once that \(\mathrm{I}(\mathbf{K}) < +\infty\) for every compact set \(\mathbf{K} \subset \mathbf{T}\) .
\item
  One can show that if I is additive, that is, satisfies (4) for finite sequences, and if I is inner regular, then I is countably additive (Exer. 7). The reader can also ascertain that only additivity and inner regularity are used in the proof of Th. 2 below.
\end{enumerate}

THEOREM 2. --- Let T be a topological space, and let I be a function defined on \(\mathfrak{B}(T)\) , with values in \(\overline{\mathbf{R}}_{+}\) . In order that there exist a measure \(\mu\) on T such that \(\mu^{\bullet}(\mathbf{A}) = \mathrm{I}(\mathbf{A})\) for every \(\mathbf{A} \in \mathfrak{B}(T)\) , it is necessary and sufficient that I be countably additive, locally bounded and inner regular. The measure \(\mu\) is then unique.

These three conditions are necessary: for, the mapping \(A \mapsto \mu^{\bullet}(A)\) on \(\mathfrak{B}(T)\) is countably additive ( \(\S1\) , No.~5, Cor. of Prop. 4), locally bounded by the definition of measures ( \(\S1\) , No.~2, Def. 5), and inner regular by Remark 3 of \(\S1\) , No.~2.

We pass to existence. It is clear that the restriction of I to \(\mathcal{K}(\mathrm{T})\) satisfies conditions 1), 2), 3) and 5) of the statement of Th. 1; let us show that 4) is satisfied as well. Let K be a compact subset of T, the intersection of a decreasing directed family \((\mathrm{K}_{\alpha})_{\alpha \in \mathrm{A}}\) of compact sets, and let \(\varepsilon\) be a number \(>0\) ; I being locally bounded, there exists an open (hence Borel) neighborhood V of K such that \(\mathrm{I(V)} < +\infty\) , and then there exists an index \(\alpha\) such that \(\mathrm{K}_{\alpha} \subset \mathrm{V}\) ; changing notation if necessary, we can suppose that \(\mathrm{K}_{\alpha} \subset \mathrm{V}\) for all \(\alpha \in \mathrm{A}\) . By the inner regularity of I, there exists a compact set \(\mathrm{L} \subset \mathrm{V} - \mathrm{K}\) such that \(\mathrm{I(L)} \geqslant \mathrm{I(V} - \mathrm{K)} - \varepsilon\) ; since L does not intersect K, there exists an index \(\alpha\) such that \(\mathrm{L} \cap \mathrm{K}_{\alpha} = \varnothing\) , and one then has \(\mathrm{I}(\mathrm{V} - \mathrm{K}_{\alpha}) \geqslant \mathrm{I}(\mathrm{L}) \geqslant \mathrm{I}(\mathrm{V} - \mathrm{K}) - \varepsilon\) . Since \(\mathbf{K}_{\alpha} \subset \mathbf{V}\) , it follows that \(\mathrm{I}(\mathbf{K}_{\alpha}) \leqslant \mathrm{I}(\mathbf{K}) + \varepsilon\) and the condition 4) is verified.

By Th. 1, there exists a measure \(\mu\) such that \(\mu^{\bullet}(\mathrm{K}) = \mathrm{I}(\mathrm{K})\) for all \(\mathbf{K} \in \mathfrak{K}(\mathbf{T})\) . The inner regularity of the set functions \(\mu^{\bullet}\) and I on \(\mathfrak{B}(\mathbf{T})\) then implies that \(\mu^{\bullet}(\mathbf{A}) = \mathrm{I}(\mathbf{A})\) for all \(\mathbf{A} \in \mathfrak{B}(\mathbf{T})\) , and existence is proved. The uniqueness of \(\mu\) follows from the uniqueness assertion of Th. 1.

